%\section{Parameter Groups and Adjoint Groups.}
On comparing (8) with (7), we see that a close formal analogy exists between
the functions \(\mu\) and \(u_0\) In fact, the \(\varphi(a, b)\) of (2) which
connect the parameters according to the composition law of the group may
themselves be considered as defining a group in the same way as does (1'),
\[
{a^{\prime}}^{\rho}=\varphi^{\rho}(a ; b)
\]
This relation can be regarded as a mapping of the a onto the a' according to a
transformation whose parameter is b. We shall prove that these transformations
form a group \(P_1\), which is isomorphic with \(S\) and is called the \underline{first}
\underline{parameter} \underline{group}. Indeed, if \(a=\varphi(a ; b)\) and \(a^{\prime
\prime}=\varphi(a ; ; c)\) then
\[
a^{\prime \prime}=\varphi(\varphi(a ; b) ; c)=\varphi(a ; \varphi(b ; c))
\]
where the last equality follows from the associative property of the
transformations \(f^i\). We thus see that the law of composition is the same
for the first parameter group and the original group of transformations \(f^i\).

The analogous group of transformations on the argument \(b\) of \(\varphi(a ;
b)\) is called the \underline{second} \underline{parameter} \underline{group}
second paramete\(P_2\). \(P_2\) is anti-isomorphic with \(S\),
that is, it is isomorphic when the factors are taken in the reverse
order. But since \((x y)^{-1}=y^{-1} x^{-1}\) and since a group contains
\(x^{-1}\) if it contains \(x\), the two are in fact isomorphic. Let a (or c)
be a transformation belonging to the first (or second) parameter group. Let the
operation of \(P_1\) transform \(b\) into \(b:=\varphi(b, a)\) and let the
operation of \(P_2\) transform \(b^{\prime}\) into
\(b^{\prime \prime}=\varphi\left(c, b^{\prime}\right)\). Then it is clear that
\[
b^{\prime \prime}=\varphi\left(c, b^{\prime}\right)=
\varphi(c, \varphi(b, a))=\varphi(\varphi(c, b), a),
\]
hence every element of \(P_1\) commutes with every element of \(P_2\).
The \(\mu^{\prime} s\) are the velocity field of \(P_1\) (see equation (\ref{eq:8})),
and they define the infinitesimal operator
\begin{equation}\label{eq:15}
A_\tau=\mu_\tau^{\rho} (a) \frac{\partial}{\partial a^{\rho}} .
\end{equation}
Correspondingly, in \(\mathrm{P}_2\)
\begin{equation}\label{eq:16}
B_\tau={\overline{\mu}}^{\rho}_{\tau} (b) \frac{\partial}{\partial b^{\rho}}
\end{equation}

Another operation which it is useful to consider is conjugation: Given an
element \(S_a\) of \(S\), to every element \(S_b\) of the group. there
corresponds an element \(S_b=S_a S_b S_a^{-1}\). The operation
\(b \rightarrow b^{\prime}\) is a faithful mapping of the group onto itself which depends on
\(S_a\) and which is called conjugation of \(S\) by \(S_a\). Consider now the
set of conjugations obtained by letting \(S_a\) run through all the elements of
\(S\). These conjugations themselves constitute a group of transformations,
parametrized by the parameters of \(S_a\), but not in general isomorphic with
\(S\). It is easily seen that isomorphism between \(S\) and the group of
conjugations holds if and only if the identity is the only element of \(S\)
that commutes with all elements of \(S_o\).

If we regard the relation \(x^{\prime}=S_a x\) as a coordinate transformation,
it is well known that the effect of operating with \(S_b\) on a function of
\(x^{\prime}\) is given in terms of \(x\) by the operation of the conjugate
of \(S_b\) by \(S_a\) acting on the same function of \(x\) :
\[
S_b x^{\prime}=S_a S_b x=\left(S_b x\right)^{\prime}.
\]
The conjugation gives the change in the parameters of an operation if this
operation is considered in the new system of coordinates \(x^8\). The advantage
of the group of conjugations over the parameter group is that the conjugate
element \(S_a S_e S_a^{-1}\) is infinitesimal if \(S_e\) is infinitesimal,
irrespective of the magnitude of \(S_a\).

If in the first system of coordinates \(S_e\) is expressed by
\(S_e=1+\varepsilon e^\rho X_\rho\), then, after the transformation with
\(S_a\) the same transformation \(S_e\) will be expressed by
\(1+\varepsilon_6{ }^{\prime} \rho_\rho^{\prime}\), and hence
\begin{equation}\label{eq:17}
e^{\prime} P_{X_p}^{\prime}=e^{p_{X_p}}
\end{equation}

The group of transformations \(e^\rho \rightarrow e^{\prime \rho}\) is called the
\underline{adjoint} \underline{group}. We wish to determine its infinitesimal
operators, produced by transformations \(S_a\) in the neighborhood of the identity. With
\(S_a=1+\delta a^{\sigma} X_\sigma\) and \(S_\rho=1+\varepsilon X_\rho\) we have
\[
S_\rho^{\prime}=1+\varepsilon X_\rho^{\prime}=
S_a S_\rho S_a^{-1}=
\left(S_a S_\rho S_a^{-1} S_\rho^{-1}\right) S_\rho=
\left(1+\left[\delta a^{\sigma} X_{\sigma}, \varepsilon X_\rho\right]\right)
\left(1+\varepsilon X_\rho\right).
\]
or
\[
X_\rho^{\prime}-X_\rho \equiv \delta X_\rho=
\delta a^{\sigma} \left[X_\sigma, X_\rho\right]=
c_{\sigma \rho}^{\tau} \delta a^{\sigma} X_\tau .
\]
by (\ref{eq:14}). From (\ref{eq:17}), we have
\[
d e^\tau X_\tau=-e^\rho d X_\rho=
e^\rho c_{\rho \sigma}^{\tau} \delta a{ }^\sigma X_\tau
\]
or
\begin{equation}\label{eq:18}
d e^\tau=e^\rho c_{\rho \sigma}^\tau \delta a^\sigma.
\end{equation}
If \(\mathrm{E}_\sigma\) are the infinitesimal operators of the adjoint group,
we find by comparison of (\ref{eq:18}) with (\ref{eq:7}) and (\ref{eq:13}) that
\begin{equation}\label{eq:19}
E_{\sigma}=e^{\rho} c_{\rho \sigma}^{\tau} \frac{\partial}{\partial e^{\tau}}.
\end{equation}
